\chapter{结论与未来工作}
\label{ch:conclusion}

长上下文 LLM 服务越来越受到 KV 缓存存储、搬移与同步成本的约束。本文分别从单节点内存层级和分布式 CXL 内存池两个层面研究了这一问题。

\section{研究总结}

\textbf{LKC-CXL-PIM} 通过将注意力计算移动到存储 KV 缓存数据附近，并在非线性阶段保持整数导向的数据通路，缓解了单节点瓶颈。iNLU 以紧凑的定点近似替代浮点 Softmax 模块，而异常值感知路由路径则为少量大幅值激活保留精度。在本文评估设定下，这一组合在 128K 上下文长度下将读取延迟降低了 98.9\% 以上，并实现了 17\% 的整体逻辑层面积缩减。

\textbf{DisaggKV} 则将同样的设计思想扩展到分布式系统。通过将局部性感知页面放置与 CXL 附加 PIM 节点之间的点对点标量归约结合起来，它减少了分布式 Softmax 评估中面向主机的互联流量需求。在本文评估的多节点场景中，DisaggKV 可在大约 2.9K requests/s 之前维持 50\,ms 的尾延迟目标，在 8 个节点范围内保持接近线性的扩展趋势，并在故障注入实验中实现亚秒级恢复。

\section{未来展望}

本文仍有若干值得继续推进的方向。首先，目前的评估采用了周期级内存仿真、事件驱动 Fabric 仿真和综合代理估计相结合的混合方法；未来可以通过更完整的 RTL 或硬件原型平台，对相同思想进行进一步验证。其次，分布式调度器仍可扩展到更大规模 Fabric、更丰富的拓扑，以及更动态的工作负载组合。第三，若能补充端到端模型质量研究，例如 iNLU 近似下的 perplexity 或服务质量测量，将进一步加强体系结构效率与应用层行为之间的联系。

尽管这些问题仍待深入，本文结果已经表明，长上下文服务并不必然受限于以主机为中心的 KV 缓存搬移。整数近内存执行与基于 CXL 的点对点归约，为长上下文大语言模型系统的扩展提供了一条具有现实意义的体系结构方向。
